\documentclass[11pt]{article}

\usepackage[margin=1in]{geometry}
\usepackage{amsmath,amssymb,amsfonts}
\usepackage{booktabs}
\usepackage{enumitem}
\usepackage{hyperref}
\usepackage{microtype}
\usepackage{xcolor}

\setlist[itemize]{leftmargin=*,topsep=2pt,itemsep=2pt}
\setlist[enumerate]{leftmargin=*,topsep=2pt,itemsep=3pt}

\title{Analytical Simplifications for the Production PNP--BV System}
\author{Jake Weinstein}
\date{May 3, 2026}

\newcommand{\OO}{\mathrm{O}_2}
\newcommand{\PP}{\mathrm{H}_2\mathrm{O}_2}
\newcommand{\HH}{\mathrm{H}^{+}}
\newcommand{\ClO}{\mathrm{ClO}_4^{-}}
\newcommand{\VT}{V_T}
\newcommand{\Eeq}{E_{\mathrm{eq}}}
\newcommand{\dd}{\,\mathrm{d}}
\newcommand{\epshat}{\varepsilon}

\begin{document}
\maketitle

\begin{abstract}
The current production forward model solves a three-species log-concentration
Poisson--Nernst--Planck system with log-rate Butler--Volmer boundary kinetics
and an analytical Boltzmann \(\ClO\) counterion. A closed-form solution of the
full two-dimensional steric PNP--BV system is not realistic, but the model
does contain several exploitable analytical structures. The main simplification
is a matched-asymptotic reduction: represent the diffuse Debye layer with an
analytical Poisson--Boltzmann profile, solve only an outer electroneutral
diffusion--reaction problem, and initialize the full solver with the resulting
composite state. In the electroneutral outer region, the existing Boltzmann
\(\ClO\) constraint makes proton transport ambipolar:
\[
  J_{\HH} = -2D_{\HH}\nabla c_{\HH}.
\]
This collapses the one-dimensional steady outer problem to linear profiles
coupled to a two-rate Butler--Volmer algebraic system. The resulting
voltage-specific initial condition directly targets the positive-voltage
Newton failure mode: the solver must start close to the exponentially depleted
Boltzmann proton layer rather than discover it from a linear-potential guess.
\end{abstract}

\section{Production model and question}

The active production stack is the formulation used by
\texttt{Forward/bv\_solver/forms\_logc.py} and
\texttt{Forward/bv\_solver/boltzmann.py}. The dynamic species are
\[
  [\OO,\PP,\HH],
  \qquad
  z = [0,0,+1],
\]
while inert \(\ClO\) is removed as a Nernst--Planck unknown and enters
Poisson's equation analytically as
\begin{equation}
  c_-(x) = c_-^b \exp(\phi(x)),
  \label{eq:clo-boltzmann}
\end{equation}
where \(\phi=\Phi/\VT\) is the nondimensional electrostatic potential.
The remaining concentrations are represented by
\[
  u_i = \log c_i,
  \qquad
  c_i = \exp(u_i).
\]
The electrode reactions are
\begin{align}
  R_1:\quad & \OO + 2\HH + 2e^- \rightleftharpoons \PP,
  &
  \Eeq^{(1)} &= 0.68~\mathrm{V~vs~RHE}, \\
  R_2:\quad & \PP + 2\HH + 2e^- \rightarrow 2\mathrm{H}_2\mathrm{O},
  &
  \Eeq^{(2)} &= 1.78~\mathrm{V~vs~RHE}.
\end{align}
The question is whether this PNP--BV system admits analytical simplifications
similar to the Boltzmann counterion reduction, whether the whole system can be
solved analytically, and whether a useful analytical initial condition can be
constructed for a specified voltage and experimental setup.

\section{Main conclusion}

A closed-form solution of the full production system should not be expected.
The full problem is nonlinear, singularly perturbed, steric, and
two-dimensional, with finite reaction fluxes coupled through nonlinear
Butler--Volmer boundary rates. Even a one-dimensional steady PNP--BV problem
with nonzero reactive flux generally does not reduce to elementary closed
form.

The useful analytical object is instead a matched-asymptotic composite state:
\begin{enumerate}
  \item Use Poisson--Boltzmann or Gouy--Chapman theory for the charged Debye
        layer adjacent to the electrode.
  \item Use an electroneutral outer problem away from that layer.
  \item In the outer region, use the analytic \(\ClO\) relation to collapse
        proton electromigration into ambipolar diffusion.
  \item Couple the resulting linear one-dimensional outer profiles to a small
        Butler--Volmer algebraic system for \(R_1\) and \(R_2\).
\end{enumerate}
This is not a replacement for the full PNP--BV solve. It is a high-quality
analytical predictor and initial condition, especially at positive voltages
where the current linear initial potential profile is far from the physical
Boltzmann layer.

\section{Exact zero-flux simplification}

Ignoring steric terms for the moment, the log-concentration Nernst--Planck
flux can be written as
\begin{equation}
  J_i
  =
  -D_i c_i \nabla(u_i + z_i\phi).
\end{equation}
Define the electrochemical-potential-like variable
\begin{equation}
  \mu_i = u_i + z_i\phi.
\end{equation}
Then
\begin{equation}
  J_i = -D_i c_i \nabla\mu_i.
  \label{eq:mu-flux}
\end{equation}
Therefore any zero-flux species in a connected region satisfies
\[
  \nabla\mu_i = 0,
  \qquad
  c_i = C_i \exp(-z_i\phi).
\]
Equation~\eqref{eq:clo-boltzmann} is exactly this simplification applied to
the inert supporting electrolyte counterion. The same argument applies
locally to \(\HH\) inside the diffuse layer whenever the normal reaction flux
is small relative to rapid double-layer equilibration:
\begin{equation}
  c_{\HH}(y)
  \approx
  c_{\HH}^{\mathrm{outer}}
  \exp[-(\phi(y)-\phi_{\mathrm{outer}})].
  \label{eq:h-inner-boltzmann}
\end{equation}
This is not globally exact because protons are consumed by both reactions, but
it is precisely the approximation needed in the thin region that currently
causes Newton failures.

With steric effects active, the zero-flux potential is modified to
\[
  \mu_i = u_i + z_i\phi + \mu_{\mathrm{steric}}(c).
\]
The Boltzmann factor is then modified by an activity coefficient. For the
current dilute production cases with \(a=0.01\), the nonsteric formula is the
right first analytical initial condition. A steric Picard correction can be
added after the nonsteric profile works.

\section{Analytical Debye layer}

Let \(h_o\) be the outer proton concentration at the edge of the diffuse
layer. Electroneutrality with the analytic \(\ClO\) ion gives the outer
potential
\begin{equation}
  \phi_o = \log\!\left(\frac{h_o}{c_-^b}\right).
\end{equation}
Define the diffuse-layer potential drop
\begin{equation}
  \psi(y) = \phi(y)-\phi_o.
\end{equation}
For a symmetric monovalent \(\HH/\ClO\) layer, the local Boltzmann profiles
are
\begin{equation}
  c_{\HH} = h_o \exp(-\psi),
  \qquad
  c_- = h_o \exp(+\psi).
  \label{eq:symmetric-boltzmann}
\end{equation}
The one-dimensional inner Poisson equation becomes
\begin{equation}
  \epshat\,\psi'' = 2h_o\sinh\psi,
  \qquad
  \lambda_{D,o} = \sqrt{\frac{\epshat}{2h_o}},
  \label{eq:pb-inner}
\end{equation}
where \(\epshat\) is the nondimensional Poisson coefficient
\((\lambda_D/L)^2\) in the code.

On a semi-infinite domain, the Gouy--Chapman solution is
\begin{equation}
  \tanh\!\left(\frac{\psi(y)}{4}\right)
  =
  \tanh\!\left(\frac{\psi_D}{4}\right)
  \exp\!\left(-\frac{y}{\lambda_{D,o}}\right),
  \qquad
  \psi_D = \phi_e-\phi_o,
  \label{eq:gouy-chapman-tanh}
\end{equation}
or equivalently
\begin{equation}
  \psi(y)
  =
  4\,\operatorname{atanh}
  \left[
    \tanh\!\left(\frac{\psi_D}{4}\right)
    \exp\!\left(-\frac{y}{\lambda_{D,o}}\right)
  \right].
  \label{eq:gouy-chapman}
\end{equation}
For small voltage, the Debye--Huckel approximation is enough:
\[
  \psi(y) \approx \psi_D\exp(-y/\lambda_{D,o}),
  \qquad
  c_{\HH} \approx h_o(1-\psi),
  \qquad
  c_- \approx h_o(1+\psi).
\]
For the production voltage grid, especially on the positive side, the full
Gouy--Chapman expression is the better default.

\section{Combined Boltzmann--Tafel factor}

Both production reactions contain the cathodic proton factor
\((c_{\HH}/c_{\HH}^{\mathrm{ref}})^2\). In the diffuse-layer approximation,
\[
  c_{\HH,s} = h_o\exp(-\psi_D).
\]
Thus each cathodic Butler--Volmer branch contains
\begin{equation}
  \left(\frac{c_{\HH,s}}{c_{\HH}^{\mathrm{ref}}}\right)^2
  \exp(-\alpha_j n\eta_j)
  =
  \left(\frac{h_o}{c_{\HH}^{\mathrm{ref}}}\right)^2
  \exp(-2\psi_D-\alpha_j n\eta_j).
  \label{eq:boltzmann-tafel}
\end{equation}
This compactly explains the positive-voltage conditioning wall. At anodic
\(\psi_D\gg 1\), the cathodic branches are suppressed not only by the usual
Tafel factor, but also by two-proton Boltzmann depletion. At
\(V_{\mathrm{RHE}}=+1.0~\mathrm{V}\), \(\psi_D\) can be \(O(40)\), so a
physically reasonable initial condition has \(u_{\HH}\) tens of log-units
below bulk at the electrode.

\section{Outer electroneutral simplification}

Outside the Debye layer, electroneutrality gives
\begin{equation}
  c_{\HH} \approx c_-^b\exp(\phi),
  \qquad
  \phi = \log\!\left(\frac{c_{\HH}}{c_-^b}\right).
  \label{eq:outer-electroneutral}
\end{equation}
Insert this relation into the proton Nernst--Planck flux:
\begin{align}
  J_{\HH}
  &=
  -D_{\HH}c_{\HH}\nabla(\log c_{\HH}+\phi) \\
  &=
  -D_{\HH}c_{\HH}
  \nabla\!\left(\log c_{\HH}+\log\!\frac{c_{\HH}}{c_-^b}\right) \\
  &=
  -2D_{\HH}\nabla c_{\HH}.
  \label{eq:ambipolar-h}
\end{align}
Therefore the outer proton equation is ordinary diffusion with effective
diffusivity \(2D_{\HH}\). The neutral species are already ordinary diffusion:
\[
  J_{\OO}=-D_{\OO}\nabla c_{\OO},
  \qquad
  J_{\PP}=-D_{\PP}\nabla c_{\PP}.
\]
In a one-dimensional steady outer approximation, all three outer
concentrations are linear functions of distance from the electrode.

\section{Two-rate outer algebraic model}

Let \(y=0\) be the electrode and \(y=1\) the bulk. Let \(R_1\) and \(R_2\) be
positive in the cathodic directions used by the implementation:
\[
  R_1:\OO\rightarrow\PP,
  \qquad
  R_2:\PP\rightarrow\mathrm{H}_2\mathrm{O}.
\]
The production stoichiometries are
\[
  s_{R_1}=[-1,+1,-2],
  \qquad
  s_{R_2}=[0,-1,-2].
\]
For the linear outer profiles, the electrode-side outer concentrations are
\begin{align}
  O_s &= O_b - \frac{R_1}{D_{\OO}}, \\
  P_s &= P_b + \frac{R_1-R_2}{D_{\PP}}, \\
  h_o &= H_b - \frac{R_1+R_2}{D_{\HH}}.
  \label{eq:outer-surface}
\end{align}
The proton expression uses the ambipolar diffusivity \(2D_{\HH}\) and the
two-proton stoichiometry, so the factors cancel.

The proton concentration entering the BV rate is the outer proton value
multiplied by the inner Boltzmann factor:
\begin{equation}
  H_s = h_o\exp(-\psi_D),
  \qquad
  \psi_D = \phi_e-\phi_o,
  \qquad
  \phi_o = \log(h_o/c_-^b).
  \label{eq:h-surface}
\end{equation}
For the current no-Stern production stack, the code evaluates
\[
  \eta_j = \phi_{\mathrm{applied}}-\Eeq^{(j)}
\]
in the BV exponent. If a Stern/Frumkin formulation is enabled, this should be
replaced by that formulation's corrected \(\eta_j=\phi_m-\phi_s-\Eeq^{(j)}\).

Holding \(H_s\) fixed for one Picard step, define
\begin{align}
  A_1 &=
  k_1
  \left(\frac{H_s}{H_{\mathrm{ref}}}\right)^2
  \exp(-\alpha_1 n\eta_1),
  &
  B_1 &=
  k_1\exp((1-\alpha_1)n\eta_1), \\
  A_2 &=
  k_2
  \left(\frac{H_s}{H_{\mathrm{ref}}}\right)^2
  \exp(-\alpha_2 n\eta_2).
\end{align}
The reduced BV rates are
\[
  R_1 = A_1O_s-B_1P_s,
  \qquad
  R_2 = A_2P_s.
\]
Substituting \eqref{eq:outer-surface} gives a \(2\times 2\) linear system:
\begin{equation}
\begin{bmatrix}
1 + A_1/D_{\OO} + B_1/D_{\PP} & -B_1/D_{\PP} \\
-A_2/D_{\PP}                  & 1 + A_2/D_{\PP}
\end{bmatrix}
\begin{bmatrix} R_1 \\ R_2 \end{bmatrix}
=
\begin{bmatrix}
A_1O_b-B_1P_b \\
A_2P_b
\end{bmatrix}.
\label{eq:two-rate-system}
\end{equation}
This algebraic solve gives a voltage-specific approximation to surface rates
and concentrations. Updating \(h_o\), \(\phi_o\), \(\psi_D\), and the
coefficients \(A_1,A_2,B_1\) by fixed point yields a cheap nonlinear reduced
model.

Useful limiting cases follow immediately:
\begin{itemize}
  \item If \(R_1\) is one-way cathodic and \(H_s\) is fixed, then
        \[
          R_1 = \frac{A_1O_b}{1+A_1/D_{\OO}},
        \]
        recovering the transition from kinetic control to the \(\OO\)
        diffusion limit.
  \item If \(R_2\) is one-way and \(R_1\) is already estimated, then
        \[
          R_2 =
          \frac{A_2(P_b+R_1/D_{\PP})}{1+A_2/D_{\PP}}.
        \]
  \item If anodic \(R_1\) dominates, the same algebraic system enforces
        peroxide depletion through \(P_s\); the rate cannot exceed the
        peroxide supply implicit in \(P_s\ge 0\).
\end{itemize}

\section{Voltage-specific analytical initial condition}

The most practical output is an initial condition for the full log-c PNP--BV
solver. For a specified voltage and experiment:
\begin{enumerate}
  \item Convert voltage to nondimensional potential
        \(\phi_e=V_{\mathrm{RHE}}/\VT\).
  \item Initialize \(R_1=R_2=0\), \(h_o=H_b\), \(\phi_o=0\), and
        \(\psi_D=\phi_e\).
  \item Compute \(H_s=h_o\exp(-\psi_D)\).
  \item Build \(A_1\), \(B_1\), and \(A_2\) from the BV parameters.
  \item Solve \eqref{eq:two-rate-system} for \(R_1,R_2\).
  \item Update
        \[
          O_s=O_b-\frac{R_1}{D_{\OO}},
          \quad
          P_s=P_b+\frac{R_1-R_2}{D_{\PP}},
          \quad
          h_o=H_b-\frac{R_1+R_2}{D_{\HH}},
        \]
        and then
        \[
          \phi_o=\log(h_o/c_-^b),
          \qquad
          \psi_D=\phi_e-\phi_o.
        \]
  \item Under-relax and repeat until the rates stop changing.
  \item Construct outer linear profiles:
        \begin{align}
          O_{\mathrm{out}}(y) &= O_s+(O_b-O_s)y, \\
          P_{\mathrm{out}}(y) &= P_s+(P_b-P_s)y, \\
          H_{\mathrm{out}}(y) &= h_o+(H_b-h_o)y.
        \end{align}
  \item Construct the Debye correction:
        \[
          \lambda_{D,o}=\sqrt{\frac{\epshat}{2\max(h_o,h_{\min})}},
        \]
        \[
          \psi(y)
          =
          4\,\operatorname{atanh}
          \left[
            \tanh\!\left(\frac{\psi_D}{4}\right)
            \exp\!\left(-\frac{y}{\lambda_{D,o}}\right)
          \right].
        \]
  \item Build the composite initial condition:
        \begin{align}
          \phi_{\mathrm{IC}}(y)
          &=
          \log\!\left(\frac{H_{\mathrm{out}}(y)}{c_-^b}\right)
          + \psi(y), \\
          H_{\mathrm{IC}}(y)
          &=
          H_{\mathrm{out}}(y)\exp[-\psi(y)], \\
          O_{\mathrm{IC}}(y)
          &=
          O_{\mathrm{out}}(y), \\
          P_{\mathrm{IC}}(y)
          &=
          \max(P_{\mathrm{out}}(y),P_{\min}), \\
          u_{i,\mathrm{IC}}(y)
          &=
          \log c_{i,\mathrm{IC}}(y).
        \end{align}
\end{enumerate}
For the two-dimensional rectangle, this is used as a \(y\)-only profile and
interpolated across \(x\).

This IC places Newton near the correct Debye-layer manifold from the start.
It also respects the outer reaction-diffusion mass balances instead of
setting every concentration to its bulk value.

\section{\texorpdfstring{Can \(\HH\) be removed analytically like \(\ClO\)?}{Can H+ be removed analytically like ClO4-?}}

Not exactly in the full model. Protons participate in both BV reactions
through the \((c_{\HH}/c_{\HH}^{\mathrm{ref}})^2\) factors and have nonzero
reactive boundary flux. Removing them globally would erase a real transport
limitation.

There are three controlled reductions:
\begin{enumerate}
  \item \textbf{Inner-layer Boltzmann \(\HH\).} Use the proton Boltzmann
        relation only inside the Debye layer. This is the least invasive and
        most useful initialization strategy.
  \item \textbf{Outer ambipolar \(\HH\).} Replace outer proton NP transport
        with \(J_{\HH}=-2D_{\HH}\nabla c_{\HH}\). This keeps proton depletion
        while removing outer Poisson stiffness.
  \item \textbf{All-Boltzmann \(\HH\) surrogate.} Set
        \(H_s=H_b\exp(-\psi_D)\) directly and drop proton transport. This may
        be a fast surrogate or diagnostic, but it is a different physical
        model.
\end{enumerate}
The recommended model-development path is the first two together: an outer
electroneutral diffusion--BV solve plus analytical Debye-layer matching.

\section{Relationship to Stern-layer support}

The current no-Stern production mode imposes Dirichlet
\(\phi=\phi_{\mathrm{applied}}\) at the electrode. This puts the full applied
potential into the resolved diffuse layer, making large positive voltages
extremely stiff.

Physically, this is a strong assumption. It identifies the solution-side
potential at the reaction plane with the metal potential:
\[
  \phi_s = \phi_m = \phi_{\mathrm{applied}}.
\]
Real electrode/electrolyte interfaces normally contain a compact Stern layer.
Ions have a finite closest-approach distance, solvent is structured near the
surface, and some voltage can drop across this compact region before the
diffuse PNP layer begins. The applied metal-to-bulk voltage is therefore more
naturally split as
\[
  \phi_m-\phi_{\mathrm{bulk}}
  =
  \Delta\phi_{\mathrm{Stern}}
  +
  \Delta\phi_{\mathrm{Diffuse}}.
\]
The current Dirichlet model effectively sets
\(\Delta\phi_{\mathrm{Stern}}=0\), so the resolved diffuse layer must absorb
the whole interfacial potential drop.

A Stern formulation replaces that electrode Dirichlet condition with a
capacitance relation and evaluates the BV overpotential as
\[
  \eta = \phi_m-\phi_s-\Eeq.
\]
The capacitance law is
\[
  \sigma = C_S(\phi_m-\phi_s),
\]
where \(C_S\) is the Stern-layer capacitance. In the PNP weak form this is a
Robin-type electrostatic boundary condition, not a prescribed-potential
condition. The comparison is:
\[
\begin{array}{ll}
\text{Dirichlet:} & \text{prescribe the solution potential } \phi_s, \\
\text{Stern:}     & \text{prescribe the metal potential } \phi_m
                    \text{ and relate diffuse-layer charge to }
                    \phi_m-\phi_s.
\end{array}
\]
This is generally more physical for an electrochemical interface, but it
introduces the additional physical parameter \(C_S\). If \(C_S\) is poorly
chosen, the model can be structurally more physical while still being less
predictive.

The limiting behavior is important:
\[
  C_S\rightarrow\infty
  \quad\Longrightarrow\quad
  \phi_m-\phi_s\rightarrow 0,
\]
which recovers the current Dirichlet-like no-Stern behavior. For finite
\(C_S\), \(\phi_s\) is not pinned to \(\phi_m\); some applied voltage can drop
in the compact layer. For the positive-voltage solver problem, this matters
because it can keep the resolved solution-side potential \(\phi_s\) much
smaller than the metal voltage \(\phi_m\), reducing the extreme Boltzmann
source \(c_-^b\exp(\phi_s)\) that currently drives the Poisson residual wall.

The same analytical framework applies, but \(\psi_D\) is no longer simply
\(\phi_e-\phi_o\). It should be solved with a compact capacitance relation:
\[
  \sigma_{\mathrm{GC}}(\psi_D)
  =
  C_S(\phi_m-\phi_s),
  \qquad
  \phi_s=\phi_o+\psi_D.
\]
The Gouy--Chapman surface charge follows from the first integral,
\[
  \sigma_{\mathrm{GC}}
  \sim
  \operatorname{sign}(\psi_D)\sqrt{8\epshat h_o}
  \sinh\!\left(\frac{|\psi_D|}{2}\right),
\]
up to the exact sign convention used by the weak Poisson boundary term. This
is more than an initializer: it is a route to replacing the explicit Debye
layer with an algebraic boundary condition.

\section{What is solved vs.\ what is analytical}

The safest first use is not to replace the model. The analytical
Debye/Stern construction should first be used as an initial condition for the
existing full PNP--BV solver.

In the current resolved model, the PDE solver computes the outer transport,
the explicit Debye layer, and the BV surface flux. In the proposed
analytical-initial-condition workflow, the analytical formula computes only a
better starting profile for the Debye-layer potential and proton correction;
the same full PNP--BV equations are still solved afterward.

A later reduced model would make a stronger split:
\[
\begin{array}{ll}
\text{PDE solver:}
  & \text{outer } \OO,\PP,\HH \text{ transport, and possibly an outer} \\
  & \text{electroneutral potential relation}, \\
\text{Analytical boundary map:}
  & \phi_s \text{ from Stern/diffuse-layer charge balance}, \\
  & H_s = H_o\exp[-(\phi_s-\phi_o)], \\
  & O_s\approx O_o,\quad P_s\approx P_o, \\
\text{BV flux law:}
  & \OO\text{ flux}=-R_1, \\
  & \PP\text{ flux}=R_1-R_2, \\
  & \HH\text{ flux}=-2R_1-2R_2.
\end{array}
\]
Thus the algebraic Debye/Stern relation does not replace the Butler--Volmer
boundary condition. It only replaces, or initially approximates, the thin
electrostatic layer that maps outer electrolyte values to reaction-plane
values. BV remains the reactive flux law.

This distinction is important. As an initial condition, the analytical Debye
profile cannot change the physical solution if the full solver converges; it
only improves basin access. As a reduced boundary model, it becomes a modeling
approximation and must be validated against resolved PNP--BV wherever the
resolved solver is available.

\section{Implementation recommendation}

The lowest-risk next step is an optional \texttt{debye\_boltzmann}
initializer for the log-c stack:
\begin{enumerate}
  \item Use the current mesh coordinate \(y\), nondimensional scaling, BV
        reaction config, and Boltzmann counterion config.
  \item Start with the no-Stern formula above and ignore steric in the IC.
  \item Floor \(O_s,P_s,h_o,H_s\) only before logarithms and report any floor
        hit, since that identifies a diffusion-limit regime.
  \item Test first near and above \(+0.60~\mathrm{V}\), where the current
        linear IC is far from the Boltzmann layer.
  \item Compare against already converged voltages by measuring SNES
        iterations and final observables. The IC should not change the final
        converged solution; it should only improve basin access.
\end{enumerate}

\section{Answers to the motivating questions}

\paragraph{Can the whole system be solved analytically?}
No. Not for the full two-dimensional production PNP--BV model with finite
reaction fluxes, steric terms, and two coupled nonlinear BV reactions.

\paragraph{Can part of the system be solved analytically?}
Yes. The zero-flux charged layer has a Boltzmann form; the symmetric
one-dimensional Debye layer has a Gouy--Chapman solution; the outer
electroneutral proton equation reduces to ambipolar diffusion; and the
one-dimensional steady outer reaction-diffusion problem reduces to a
two-rate algebraic system.

\paragraph{Can we get a good analytical IC for a given voltage?}
Yes. The matched Debye-layer plus outer algebraic profile is likely the
highest-value analytical simplification for the current solver. It directly
targets the failure mode documented in
\texttt{docs/peroxide\_window\_investigation.md}: Newton starts too far from
the exponentially depleted proton layer and the exponentially enriched
counterion layer. The proposed IC starts close to that manifold.

\end{document}
